\chapter{The master dictionary}
\index{master dictionary}\index{comparison!non-Archimedean frameworks}
\label{chap:master-dictionary}

The word ``dictionary'' is dangerous in pluripotential theory.
It can mean an equivalence of categories, a construction in one direction,
or only a useful resemblance between formulas.
This chapter separates those three meanings.
The resulting diagram is deliberately smaller than the collection of all
objects introduced in the book.
Every solid identification below is made on a locus where both sides are
defined and where a published comparison theorem applies.

Four warnings govern the chapter.
First, the BFJ theory over a nontrivially valued field is not obtained by
putting the absolute value equal to one in its formulas.
Second, the CLD differential calculus does not presently give a global
equivalence with every singular BFJ-psh metric.
Third, the common underlying valuation space in projective complex geometry
does not by itself identify the DXZ and Mesquita--Piccione psh cones.
Fourth, the normalized BFJ local link and the Reeb-normalized affine-cone
link cut different radial actions and may not be identified by notation.

\section{Arrow types and domain signatures}

We use the following three arrow styles throughout:
\[
 \mathcal A\xleftrightarrow{\ \sim\ }\mathcal B,
 \qquad
 \mathcal A\longrightarrow\mathcal B,
 \qquad
 \mathcal A\dashrightarrow\mathcal B.
\tag{63.1}\label{eq:dict-arrow-legend}
\]
The first is a proved equivalence or bijection under the displayed
hypotheses.
The second is a proved construction in the indicated direction, with no
inverse asserted.
The dashed arrow means analogy only.
In particular, a dashed arrow may not be used to transport a theorem.

We now name the domains that occur in the diagram.

\begin{definition}[Nontrivially valued model domain]
\label{def:dict-ntv-model-domain}
Let \(\mathsf D_{\mathrm{mod}}(K)\) consist of tuples
\[
 (K,X,L_1,\ldots,L_n,\mathcal X,
       \mathcal L_1,\ldots,\mathcal L_n)
\]
such that \(K\) is complete and discretely valued,
\(X\) is a proper reduced \(n\)-dimensional \(K\)-scheme,
\(\mathcal X\) is a proper flat \(K^\circ\)-model relatively normal in
\(X\), and \(\mathcal L_i\) extends \(L_i^{\otimes m_i}\) for some
\(m_i>0\).
The model metric on \(L_i\) is the rational metric obtained by dividing the
formal metric of \(\mathcal L_i\) by \(m_i\).
\end{definition}

Relative normality is included because it gives one divisorial analytic
point above the generic point of each special-fibre component.
The more general CLD formula allows nonreduced formal models and inserts
defect factors; the clean Chambert--Loir formula uses the domain just stated.

\begin{definition}[CLD formal domain]
\label{def:dict-cld-domain}
Let \(\mathsf D_{\mathrm{CLD}}(K)\) consist of a flat formal
\(K^\circ\)-scheme \(\mathfrak X\), locally topologically of finite
presentation and pure of relative dimension \(n\), whose generic fibre is
boundaryless, together with strictly integral PL formal line bundles
\(\mathfrak L_1,\ldots,\mathfrak L_n\).
\end{definition}

Proper algebraic models in
\(\mathsf D_{\mathrm{mod}}(K)\) enter
\(\mathsf D_{\mathrm{CLD}}(K)\) by completion along the special fibre.
This is a functor, not an equivalence between all algebraic and formal
models.

\begin{definition}[Original BFJ domain]
\label{def:dict-bfj-domain}
The signature \(\mathsf D_{\mathrm{BFJ}}\) means:
\(K\) is complete and discretely valued, the residue field has
characteristic zero, \(X\) is smooth and projective over \(K\), and
\(L\) is ample.
A semipositive closed model form \(\theta\) with generic class
\(c_1(L)\) is fixed.
\end{definition}

These are the hypotheses of the foundational singular theory used in
\Crefrange{chap:bfj-model-forms}{chap:bfj-continuity}.
Later extensions must be cited separately and are not silently folded into
the signature.

\begin{definition}[Trivially valued projective domain]
\label{def:dict-bj-domain}
The signature \(\mathsf D_{\mathrm{BJ}}\) means that \(k\) is trivially
valued, \(X\) is an integral projective \(k\)-variety of pure dimension
\(n\), and \(L\) is ample or, where explicitly stated, a
\(\mathbf Q\)-line bundle in a specified numerical class.
Test configurations carry a fixed punctured-product trivialization and a
fixed linearization.
\end{definition}

Fixing the linearization is essential.
Twisting it by the character of weight \(c\) translates the corresponding
metric by \(c\).

\begin{definition}[Projective complex comparison domain]
\label{def:dict-complex-domain}
The signature \(\mathsf D_{\mathbf C}^{\mathrm{proj}}\) means that
\(X\) is a connected projective manifold over \(\mathbf C\),
\(L\) is a pseudo-effective \(\mathbf Q\)-line bundle,
\(h\) is smooth on \(L\), and \(\theta=c_1(L,h)\).
\end{definition}

This is the domain of the Darvas--Xia--Zhang comparison used in
\Cref{chap:dxz-bj-comparison}.
It does not include a singular projective variety, a general compact
K\"ahler manifold, or an arbitrary transcendental class.

\begin{definition}[Tropical-spectrum domain]
\label{def:dict-mp-domain}
The signature \(\mathsf D_{\mathrm{trop}}\) means a reduced irreducible
compact complex space equipped with the idempotent semiring of coherent
ideals and its tropical spectrum \(X^\bethan\).
\end{definition}

No positivity is part of this signature.
It supplies semivaluations, supports, centers, topology, and functoriality.

\section{How to read a comparison arrow}
\label{sec:dict-reading-arrows}

A comparison statement has six coordinates.  Suppressing any one of them is
the most common way to turn a correct theorem into a false slogan.

\begin{enumerate}
\item The \emph{base norm} says whether the ground field is nontrivially or
trivially valued.  It fixes the radial normalization of divisorial
valuations.

\item The \emph{carrier} says whether the object lives on a Berkovich space,
on an inverse limit of models, on a tropical spectrum, or on a hybrid total
space.  A homeomorphism of carriers is not yet a comparison of positive
objects.

\item The \emph{algebraic representative} records the model, test
configuration, filtration, coherent ideal, or test curve from which the
object is built.  Its equivalence relation---domination, integral closure,
Veronese change, or equality of singularity type---is part of the datum.

\item The \emph{positivity certificate} may be relative nefness, positivity
of a CLD or delta-current, a curve test, a BFJ retraction inequality, or
generalized FS approximation.  These certificates are compared only by the
theorems cited below.

\item The \emph{limiting procedure} distinguishes uniform model
approximation, decreasing psh approximation, finite-energy truncation, and
hybrid weak convergence.  Agreement at the first stage does not identify
the later stages without a uniqueness theorem.

\item The \emph{normalization} fixes additive constants, total mass, the
linearization, and whether a complex potential is a norm weight or a
squared-norm weight.
\end{enumerate}

Thus an arrow in this chapter should be read as a six-entry packet.  For
example, the first row of \eqref{eq:dict-master-diagram} has the packet
\[
\begin{array}{c|c|c|c|c|c}
\text{base}&\text{carrier}&\text{representative}&\text{positivity}
 &\text{limit}&\text{mass}\\
\hline
K,\ |\pi|<1&X^{\An}&(\mathcal X,m_i^{-1}\mathcal L_i)
 &\mathcal L_i\text{ relatively nef}&\text{model level}
 &(L_1\cdots L_n).
\end{array}
\tag{63.2}\label{eq:dict-six-coordinate-packet}
\]
The table is not decorative: changing the base to the trivial norm replaces
the special-fibre normalization by radial homogeneity, while changing the
limit from model level to finite energy requires a separate monotone
extension theorem.

There are also three distinct uses of the word \emph{valuation} in the
comparison literature.  If
\(\mathcal X_s=\sum_Eb_EE\), the BFJ/Chambert--Loir divisorial point is
\[
                       x_E=b_E^{-1}\ord_E,
              \qquad x_E(\pi)=1.
\tag{63.3}\label{eq:dict-nontrivial-divisorial-normalization}
\]
For a test configuration over a trivially valued field, the restriction of
\(b_E^{-1}\ord_E\) to the function field of \(X\) is instead placed on a
radial cone and is moved by \(v\mapsto tv\).  In the hybrid chapter, the
book's additive valuation is
\((-\log r)b_E^{-1}\ord_E\).  These are compatible rescalings inside their
own constructions, not literally the same normalized point.

\section{The diagram with hypotheses attached}

The master diagram is a family of diagrams, not one rectangle in a single
category.
Its proved part is the following.

\begin{theorem}\label{thm:dict-master-diagram}
The following rows commute on the stated domains.
Every \(\sim\) is an equivalence after imposing the quotient written below
it; every unmarked arrow is one-way.
\[
\begin{array}{c@{\quad}c@{\quad}c}
\multicolumn{3}{c}{\mathsf D_{\mathrm{mod}}(K)}\\[2mm]
\left.
 \begin{array}{c}
 \text{vertical rational Cartier data}\\
 \text{modulo common domination}
 \end{array}
\right\}
&\xleftrightarrow{\ \sim\ }&
\left\{
 \begin{array}{c}
 \text{rational model metrics}\\
 \text{with fixed generic metric datum}
 \end{array}
\right.
\\[-1mm]
&&\ \ \searrow\ c_1(\overline L_1)\wedge\cdots\wedge
                 c_1(\overline L_n)
\\[-1mm]
&&\qquad \mathcal M_{\mathrm{CLD}}
       =\mathcal M_{\mathrm{CL}}
       =\mathcal M_{\mathrm{int}}
\\[5mm]
\multicolumn{3}{c}{\mathsf D_{\mathrm{BJ}}}\\[2mm]
\left.
 \begin{array}{c}
 \text{integrally closed test configurations}\\
 \text{modulo equivariant domination}
 \end{array}
\right\}
&\xleftrightarrow{\ \sim\ }&
\left\{
 \begin{array}{c}
 \text{rational PL/FS metrics}\\
 \text{with fixed linearization}
 \end{array}
\right.
\\[5mm]
\multicolumn{3}{c}{\mathsf D_{\mathbf C}^{\mathrm{proj}}}\\[2mm]
\PSH^{\NA}_{\mathrm{tc}}(X,\theta)
&\xleftrightarrow{\ \mathfrak A\ \sim\ }&
\PSH(L^{\An})
\\[5mm]
\multicolumn{3}{c}{X\text{ projective over }\mathbf C}\\[2mm]
X^\bethan
&\xleftrightarrow{\ \sim\ }&
X^{\mathrm{an,triv}}.
\end{array}
\tag{63.4}\label{eq:dict-master-diagram}
\]
In the first row, \(\mathcal M_{\mathrm{int}}\) is the special-fibre
intersection measure, \(\mathcal M_{\mathrm{CL}}\) the
Chambert--Loir measure, and \(\mathcal M_{\mathrm{CLD}}\) the CLD current
product.
The equality is asserted only after the model datum has entered the CLD
formal domain.
\end{theorem}

\begin{proof}
For the first equivalence, a vertical rational Cartier divisor changes a
fixed model metric by its model function.
Pullback to a dominating model changes neither metric.
Conversely, the quotient of two rational model metrics on the same generic
line bundle is represented, on a common model, by a vertical rational
Cartier divisor.

Complete the model along its special fibre.
\Cref{thm:cld-recovers-cl} identifies the CLD current product with
\[
 \frac{1}{m_1\cdots m_n}
 \sum_Y b_Y
   \bigl(c_1(\mathcal L_1|_Y)\cdots
         c_1(\mathcal L_n|_Y)\bigr)\delta_{x_Y}.
\tag{63.5}\label{eq:dict-model-measure}
\]
This is both the Chambert--Loir definition and the model intersection
measure.
Thus the triangular part of the first row commutes.

In the trivially valued row, a test configuration gives the vertical
Cartier datum of a PL metric.
Integral closure does not change its valuative function.
With the punctured-product trivialization and linearization fixed, values at
central-fibre divisorial valuations recover that divisor on a common
domination.
This is the injectivity statement of
\cite[Theorem~2.31(i)]{BJ22}; the Rees construction supplies the inverse on
rational FS data.

For the complex projective row, \Cref{thm:comp-dxz} proves that
analytification is an order isomorphism.
Surjectivity uses generalized FS approximation and hence the projective
envelope theory; it is not a formal consequence of the valuative formula.

Finally, \Cref{thm:ts-algebraic-recovery} identifies a tropical
character on coherent ideals with its semivaluation on principal ideals.
GAGA identifies algebraic and analytic coherent ideals on a projective
complex variety.
Properness supplies a center for every valuation.
The resulting continuous bijection between compact Hausdorff spaces is a
homeomorphism.
\end{proof}

\begin{remark}[Meaning of the last row]
\label{rem:dict-space-not-psh}
The last row compares underlying spaces and ideal evaluations only.
There is no vertical \(\sim\)-arrow in
\eqref{eq:dict-master-diagram} from the DXZ test-curve cone to a
Mesquita--Piccione psh cone.
Such an arrow would be an additional positivity theorem, not topology or
GAGA.
No unpublished DXZ--MP equivalence is used in this book.
\end{remark}

\section{The common model core}

The most robust meeting point of the global and differential theories is a
formal metric.
We spell out what is and is not forgotten there.

\begin{proposition}\label{prop:dict-model-core}
On \(\mathsf D_{\mathrm{mod}}(K)\), the assignment
\[
 \overline L\longmapsto c_1(\overline L)
\tag{63.6}\label{eq:dict-metric-to-current}
\]
is a functorial one-way construction.
Its \(n\)-fold product agrees with the Chambert--Loir intersection measure.
It is not an equivalence of metric spaces with current spaces.
\end{proposition}

\begin{proof}
Functoriality and multilinearity of the CLD first Chern current give the
construction.
The equality of products is
\eqref{eq:dict-model-measure}.
If a metric is translated by a constant, its Chern current is unchanged,
while the metric is not.
Thus even on the model core the current forgets at least the additive
normalization.
\end{proof}

The distinction is visible in energy.
If \(V=(L^n)\), then the unnormalized energy satisfies
\[
                         E(\phi+c)=E(\phi)+cV,
\tag{63.7}\label{eq:dict-energy-translation}
\]
whereas \(c_1(\overline L \rme^{-c})=c_1(\overline L)\).
Any dictionary that identifies a metric with its curvature current must
therefore either quotient by constants or retain an auxiliary normalization.

\begin{corollary}\label{cor:dict-continuous-common-locus}
Assume \(\mathsf D_{\mathrm{BFJ}}\), and let a continuous semipositive
metric be a uniform limit of semipositive model metrics.
Its BFJ Monge--Amp\`ere measure equals its Chambert--Loir measure.
If, in addition, a CLD Bedford--Taylor product is defined by the same
semipositive formal regularization, then that product is the same weak
limit.
\end{corollary}

\begin{proof}
Choose semipositive model weights \(\phi_j\to\phi\) uniformly.
All model measures agree by
\eqref{eq:dict-model-measure}.
BFJ monotone continuity, after replacing a rapidly convergent subsequence by
a decreasing translated sequence, identifies its limit with
\(\MA_\theta(\phi)\).
Chambert--Loir continuity identifies the same weak limit with the continuous
semipositive measure.
When the CLD product is constructed from this regularization, its weak
continuity gives the third equality.
\end{proof}

The last conditional clause is part of the theorem's domain.
It must not be inverted to claim that every bounded BFJ-psh function is a
local CLD-psh function or that every singular CLD product is the BFJ product.

\section{The local differential branch of the dictionary}
\label{sec:dict-local-differential-branch}

The phrase ``CLD approach'' is often used for several later enlargements.
For comparison purposes they must be separated.  On a boundaryless strict
analytic space the relevant diagram is
\[
\begin{array}{ccccc}
\text{CLD smooth forms}
 &\longrightarrow&
\begin{array}{c}
\text{Gubler--Jell--Rabinoff}\\[-1mm]
\text{weakly smooth forms}
\end{array}
 &\longrightarrow&
\text{currents}
\\[2mm]
\downarrow&&&&\uparrow
\\[-1mm]
\text{formal/PL metrics}
 &\longrightarrow&
\begin{array}{c}
\text{Gubler--K\"unnemann}\\[-1mm]
\text{delta-forms}
\end{array}
 &\longrightarrow&
\text{delta-currents}.
\end{array}
\tag{63.8}\label{eq:dict-local-differential-ladder}
\]
The upper left arrow is an enlargement of the class of charts.  The lower
row is an enlargement of the coefficients that permits polyhedral corner
terms.  Neither row is the definition of BFJ-psh functions.

At the common PL locus, however, there is a precise three-step check.  Let
\(h\) be the weight of a formal metric.

\begin{enumerate}
\item Relative nefness of the formal model is the algebraic positivity
certificate.

\item Its first Chern object is a positive Gubler--K\"unnemann
delta-current on the exact formal domain of
\Cref{thm:gk-formal-positivity-comparison}.

\item When \(K\) is nontrivially valued and the analytic space is good and
strict, \(h\) is classically psh in the Gubler--Rabinoff curve-test sense
exactly under the hypotheses of \Cref{thm:gr-pl-comparison}.
\end{enumerate}

This chain makes the model core visible simultaneously to intersection
theory, differential currents, and curve potential theory.  It does not
produce the dashed converse
\[
 \text{arbitrary curve-test psh}
       \dashrightarrow
 \text{arbitrary BFJ-psh}.
\tag{63.9}\label{eq:dict-forbidden-local-global-converse}
\]
Indeed, the CLD class is generated by smooth positivity and monotone
operations, the Gubler--Rabinoff class is tested after every scalar extension
on every analytic curve, and the BFJ class is controlled by model
retractions and nef envelopes.  The one proved general implication recorded
in this book is the CLD-to-curve-test arrow of
\Cref{thm:cld-to-classical-psh}; no converse is silently added.

\paragraph{Source boundary for the local branch.}
The CLD smooth calculus is \cite{CLD25}.  Delta-forms and their positivity
comparison are \cite{GK17,GK19}; harmonic tropicalizations and weakly smooth
forms are \cite{GJR26}; and the PL and classical curve-test comparisons are
\cite{GR25PL,GR25Classical}.  The exact hypotheses and the non-inclusions are
spelled out in \Cref{chap:delta-forms}.  The present section only assembles
those already proved arrows and imports no singular comparison theorem.

\subsection{A multiplicity calculation at the common model core}

Write \(D=\sum_Ed_EE\) for a vertical rational Cartier divisor on
\(\mathcal X\), and let \(f_D\) be its model function.  The two
normalizations in the model formula cancel:
\[
 f_D(x_E)=\frac{d_E}{b_E},
 \qquad
 \mu_{\mathcal X}(\{x_E\})
 =\frac{b_E}{m_1\cdots m_n}
   \bigl(c_1(\mathcal L_1|_E)\cdots
         c_1(\mathcal L_n|_E)\bigr).
\tag{63.10}\label{eq:dict-multiplicity-cancellation}
\]
Consequently
\[
 \int f_D\,\rd\mu_{\mathcal X}
 =\frac1{m_1\cdots m_n}
   \sum_Ed_E
   \bigl(c_1(\mathcal L_1|_E)\cdots
         c_1(\mathcal L_n|_E)\bigr),
\tag{63.11}\label{eq:dict-model-duality-expanded}
\]
which is the vertical intersection number.  Omitting \(b_E\) from the
measure or from the normalization of \(x_E\) destroys this identity.  This
is the quickest local audit of a proposed CLD--Chambert--Loir comparison.

Replacing every extension by
\(\mathcal L_i^{\otimes q_i}\), representing
\(L_i^{\otimes m_iq_i}\), multiplies the intersection on each component by
\(q_1\cdots q_n\) and the denominator by the same product.  Hence the
measure is unchanged.  This tensor-power test explains why rational model
metrics, rather than chosen integral presentations, are the correct common
objects.

\section{Inverse systems and their realizations}

The fourth row of the planning table concerns inverse systems of models.
Here the arrows are less symmetric.

\begin{proposition}\label{prop:dict-bfj-skeletal}
Under \(\mathsf D_{\mathrm{BFJ}}\), a BFJ-psh function determines compatible
upper semicontinuous restrictions to every SNC dual complex, and
\[
                  \phi=\inf_{\mathcal X}\phi\circ p_{\mathcal X}.
\tag{63.12}\label{eq:dict-bfj-inverse-system}
\]
This gives an injective map from BFJ-psh functions to compatible skeletal
data satisfying the BFJ convexity, nefness, and retraction inequalities.
Its image is characterized by those inequalities.
\end{proposition}

\begin{proof}
This is the skeletal characterization
\cite[Section~7]{BFJ16}, recalled in
\Cref{thm:bfj-skeletal-characterization}.
The retraction inequality
\(\phi\leq\phi\circ p_{\mathcal X}\) is essential.
Without it, compatible convex functions on dual complexes need not recover a
psh function on all of \(X^{\An}\).
\end{proof}

The analogous BJ statement uses decreasing generalized FS approximation on
the valuation cone.
Its input is not the same inverse system: the ground valuation is trivial,
the radial scaling action is part of the space, and proper-support
semivaluations may take the value \(-\infty\).

\begin{proposition}\label{prop:dict-projective-realization}
Under \(\mathsf D_{\mathbf C}^{\mathrm{proj}}\), a compatible DXZ inverse
system of positive-mass \(\cI\)-model test curves determines a unique BJ-psh
metric, and every BJ-psh metric arises in this way.
Order, admissible decreasing limits, and the operations established in
\Cref{chap:dxz-bj-comparison} are preserved.
\end{proposition}

\begin{proof}
The asserted bijection is \Cref{thm:comp-dxz}.
Order reflection is \Cref{prop:comp-order-reflection}, which uses
homogeneous Legendre inversion and the valuative multiplier-ideal criterion.
Surjectivity uses decreasing generalized FS approximation.
Compatibility with operations and trace is proved in the cited chapter and
requires the relevant envelope theorems exactly where recorded there.
\end{proof}

\begin{remark}[Zero-mass safeguard]
The DXZ inverse system uses K\"ahler perturbations because a
pseudo-effective class may have zero volume.
No single positive-mass test curve then represents the unperturbed object.
The order isomorphism does not license division by that zero volume, nor does
it turn the positive-mass Okounkov body formula into a zero-mass formula.
\end{remark}

\section{The projective complex branch: same carrier, different positivity}
\label{sec:dict-projective-complex-branch}

Suppose \(X\) is projective over \(\mathbf C\).  The MP tropical spectrum
and the trivially valued Berkovich analytification have the same points:
a tropical character on coherent ideals is recovered from its values on
principal ideals, and a semivaluation \(v\) gives the ideal value
\[
                    v(\mathfrak a)
       =\min\{v(f):f\text{ is a local generator of }\mathfrak a\}.
\tag{63.13}\label{eq:dict-ideal-valuation-recovery}
\]
Coherence makes the minimum local and finite generation makes it independent
of the chosen generators.  Multiplication of ideals becomes addition of
values, while a sum of ideals becomes a minimum.  This explains the
underlying-space homeomorphism in the last row of
\eqref{eq:dict-master-diagram}.

The DXZ-to-BJ arrow attaches more structure to the same carrier.  Its input
is a compatible system of positive-mass \(\cI\)-model test curves; its
output is a psh metric obtained by generalized FS approximation.  In
particular, its positivity remembers the numerical class and its envelope
theory.  Formula \eqref{eq:dict-ideal-valuation-recovery} remembers neither.
The logical relation is therefore
\[
\begin{array}{ccc}
\text{DXZ test-curve datum}&\xrightarrow{\ \mathfrak A\ }&
     \text{BJ psh metric}\\
\downarrow\text{ forget positivity}&&\downarrow\text{ forget metric}\\
\text{coherent-ideal evaluations}&=&
     \text{MP/Berkovich semivaluation carrier}.
\end{array}
\tag{63.14}\label{eq:dict-dxz-mp-forgetful-square}
\]
Only the top arrow and the bottom equality are comparison theorems.  The
vertical arrows are forgetful descriptions.  Reading the square upward
would amount to reconstructing a psh metric from a bare point, which is
impossible even before one asks whether the MP and DXZ positivity notions
agree.

\paragraph{Source boundary for the projective branch.}
The top arrow is bounded by the projective manifold and
pseudo-effective \(\mathbf Q\)-line-bundle hypotheses of
\cite{DXZ23}.  The lower row is the tropical-spectrum construction of
\cite{MP24} together with projective GAGA.  No result in either source is
used here to assert a DXZ--MP equivalence of psh cones.

\section{Three interfaces that require extra boundary data}
\label{sec:dict-boundary-data-interfaces}

Three comparison arrows developed later in this part do not fit into
\eqref{eq:dict-master-diagram}.  The first two are local but normalize
different radial actions; the third exists only on the intersection of two
positivity domains.  Keeping them separate prevents the carrier of a theory
from being confused with its psh cone.

\subsection{The normalized local link and a Gauss shell}

Let \((Y,0)\) be a smooth complex germ of dimension \(n\geq2\), let
\(\mathfrak m\) be its maximal ideal, and let
\(\mathcal V_0=\{v:v(\mathfrak m)=1\}\) be the normalized BFJ link.  After
passing to
\[
                 K=\mathbf C((t)),\qquad |t|=\rme^{-1},
\]
the Gauss construction of \Cref{lem:local-cld-gauss-core} gives a section
and retraction
\[
 \mathcal V_0\xrightarrow{\ \iota\ }
 \{\log|\mathfrak m|=-1\}
 \xrightarrow{\ p\ }\mathcal V_0,
 \qquad p\circ\iota=\operatorname{id}_{\mathcal V_0}.
\]
For every \(\mathfrak m\)-primary ideal \(\mathfrak a\), the analytic model
function \(u_{\mathfrak a}=\log\max_j|f_j|\) satisfies
\[
             u_{\mathfrak a}\circ\iota=-v(\mathfrak a).
\]
Thus the model functions agree literally; the measures agree only after the
reference ideal supplies one further boundary factor.  More precisely,
\Cref{thm:local-cld-model-comparison} identifies the BFJ mixed measure on
\(\mathcal V_0\) with the pushforward by \(p\) of the swept CLD boundary
product in \eqref{eq:local-cld-swept-product}.  The equality holds for ideal
models and algebraizable rational semipositive models, and
\Cref{prop:local-cld-bounded-closure} extends it to the explicitly defined
joint bounded closure.

Three pieces of data are indispensable: the normalization ideal
\(\mathfrak m\), the level \(-1\), and the swept reference factor.  Removing
any of them changes the carrier or the mass.  In particular, the comparison
does not identify arbitrary formal psh functions with arbitrary CLD-v2 psh
functions, and it does not select a canonical treatment of unbounded polar
mass.  The two competing continuations at a polar endpoint are exhibited in
\eqref{eq:local-cld-polar-endpoint}.

\subsection{The Reeb-polarized affine-cone branch}

Let \((X=\Spec R,T;\xi)\) be a normal polarized affine cone over a trivially
valued algebraically closed field of characteristic zero.  Wu's analytic
link and the BFJ normalized local link have visibly different signatures:
\[
 \begin{array}{ccl}
  \mathcal V_0=\{v:v(\mathfrak m)=1\}
   &:&\text{centered valuations modulo positive scaling},\\[1mm]
  X_0=\{v:v(\mathfrak m)=0\}
   &:&\text{level-set cross-section for Reeb translation}.
 \end{array}
\]
The set in the second line depends only on the pointed affine cone; after a
good torus action and a Reeb field have been fixed, it is a cross-section for
the corresponding translation orbits.  Its reference function
\(\varphi_\xi\) satisfies
\(X_0=\{\varphi_\xi=0\}\), and the model measure is
\[
 \frac{1}{\vol(\xi)}
 \dprimeddprime\varphi_1\wedge\cdots\wedge
 \dprimeddprime\varphi_{n-1}\wedge
 \dprimeddprime\max\{\varphi_\xi,0\}.
\]
This formula exhibits two solid interfaces.  First, homogeneous
Fubini--Study functions are CLD-approachable, so the displayed expression is
a literal CLD Bedford--Taylor product supported on \(X_0\).  Second, when
\(\xi\) is rational, the projective quotient identifies it with the global
normalized Chambert--Loir probability measure on the polarized orbifold
quotient, and Wu's finite-energy class and energy recover their global BJ
counterparts.  The complete domains,
normalizations, and source theorem are recorded in
\Cref{sec:local-polarized-affine-cones,thm:local-cone-finite-energy-ma}.

The resulting comparison chain is therefore
\[
 \begin{gathered}
  \mathcal H(\xi)
  \subset\mathrm{CPSH}(\xi)
  \subset\mathcal E^1(\xi),
  \\[1mm]
  \mathcal E^1(\xi)
  \xleftarrow{\ \sim\ }
  \mathcal E^1_{\mathrm{BJ}}(V_\xi,L_\xi)
  \qquad(\xi\text{ rational}).
 \end{gathered}
\]
On \(\mathcal H(\xi)\), the operator \(\MA_\xi\) is separately realized
by the normalized CLD shell product displayed above.  Thus the diagram
compares potential classes, while the preceding formula compares measures.
By contrast,
\[
 \text{BFJ formal psh on }\mathcal V_0
 \quad\nleftrightarrow\quad
 \text{affine-cone }\mathcal E^1(\xi).
\]
The non-arrow is substantive, not typographical.  No theorem identifies the
two links, transports the finite-energy class after forgetting \(\xi\), or
compares psh classes for two nonproportional Reeb fields.  The source boundary is
\cite{Wu22AffineEnergy,Wu24AffineCones}; the present subsection only places
their established arrows in the dictionary.

\subsection{The finite-energy DXZ--MP--BJ triangle}

There are two independent comparison theorems in complex geometry:
\[
\begin{array}{ccc}
 \mathcal E^{1,\NA}_{\mathrm{DXZ}}(X,\omega)
 &\xrightarrow{\ \sim\ }&
 \mathcal E^1_{\mathrm{MP}}(\alpha)\\[2mm]
 \Big\downarrow && \Big\downarrow\\[-1mm]
 \PSH^{\NA}_{\mathrm{DXZ}}(X,\theta)
 &\xrightarrow{\ \sim\ }&
 \PSH(L^{\mathrm{an,triv}}).
\end{array}
\]
The top arrow is \Cref{thm:pseff-kahler-finite-energy-link}; it is available
for a K\"ahler class on a compact K\"ahler manifold and is characterized by
rational divisorial values and equality of radial energy.  The bottom arrow
is \Cref{thm:comp-dxz}; it requires a smooth projective manifold and a
pseudo-effective \(\mathbf Q\)-line bundle, but it does not require finite
energy or positive volume.  The right vertical notation is meaningful only
when the two domains overlap; it is not a general inclusion of the MP cone
into the BJ cone.

On a smooth polarized projective manifold and in finite energy, the two
horizontal comparisons have the same rational divisorial values.  The
divisorial determination theorems on the two sides therefore identify the
corresponding objects on this common locus.  This is the common-domain
reading of the two cited comparison theorems, not an additional extension.
Outside it, the diagram is not a square of globally defined maps.  For an
ample projective class without a
finite-energy hypothesis, the full MP--BJ identification remains
\Cref{prob:pseff-projective-rational-full-link}; for a non-nef
pseudo-effective class, \Cref{prop:pseff-relative-nef-obstruction} shows
that the unchanged relatively nef MP model signature cannot even have the
required generic-fibre class.

\section{Degenerations, filtrations, and rays}

There is a second diagram inside the projective complex domain.
Its exact form is developed in \Cref{chap:rays-filtrations}; here we
record the arrow types.
\[
\begin{array}{ccccc}
\text{normal relatively ample TC}
&\longrightarrow&
\text{f.g. rational filtration}
&\longrightarrow&
\text{FS/PL metric}
\\[3mm]
&&\downarrow&&\downarrow
\\[-1mm]
&&\text{maximal test curve}
&\xleftrightarrow{\ \sim\ }&
\text{weak geodesic ray}.
\end{array}
\tag{63.15}\label{eq:dict-degeneration-diagram}
\]

The first arrow becomes an equivalence only after fixing the section ring,
taking a suitable Veronese, passing to the integrally closed Rees algebra,
and quotienting by equivariant domination.
A general linearly bounded filtration need not be finitely generated and
therefore need not come from a test configuration.
The bottom equivalence is the Ross--Witt Nystr\"om correspondence in its
stated analytic domain.

\begin{proposition}\label{prop:dict-degeneration-commutes}
Let \((X,L)\) be a polarized complex manifold and let
\((\mathcal X,\mathcal L)\) be a normal relatively ample test
configuration.
The ray obtained from its filtered section ring by maximal envelopes and
Legendre transform is the Phong--Sturm weak geodesic ray of the test
configuration.
\end{proposition}

\begin{proof}
Equivariant extension of a section gives the filtration.
Its partial Bergman kernels give the level obstacles of the test curve.
The Legendre transform of their maximal envelopes is the limit of the
weighted Bergman rays.
Ross--Witt Nystr\"om
\cite[Theorem~9.2]{RWN14} identifies this limit with the Phong--Sturm ray.
\end{proof}

This comparison is independent of the flawed invariant-theory assertion in
\cite[Theorem~5.6]{BHJ19}.
The erratum boundary is treated explicitly in
\Cref{chap:rays-filtrations}.

\pagebreak
\section{The comparison-status ledger}

The following ledger is a compact audit of the master dictionary.

\begingroup
\setlength{\parindent}{0pt}

\textbf{Pair of realizations}\par\nopagebreak
model intersection / Chambert--Loir\par\nopagebreak

\textbf{Status}\par\nopagebreak
equality\par\nopagebreak

\textbf{Exact common locus}\par\nopagebreak
rational model metrics on a proper relatively normal model over a
discretely valued field

\medskip
\textbf{Pair of realizations}\par\nopagebreak
CLD formal product / Chambert--Loir\par\nopagebreak

\textbf{Status}\par\nopagebreak
equality\par\nopagebreak

\textbf{Exact common locus}\par\nopagebreak
completion of the preceding model; more generally the CLD formal formula
has its boundaryless and defect-factor hypotheses

\medskip
\textbf{Pair of realizations}\par\nopagebreak
BFJ / Chambert--Loir\par\nopagebreak

\textbf{Status}\par\nopagebreak
equality of measures\par\nopagebreak

\textbf{Exact common locus}\par\nopagebreak
model metrics and continuous semipositive metrics obtained by uniform
semipositive model approximation

\medskip
\textbf{Pair of realizations}\par\nopagebreak
BFJ singular psh / CLD singular psh\par\nopagebreak

\textbf{Status}\par\nopagebreak
no global equivalence asserted\par\nopagebreak

\textbf{Exact common locus}\par\nopagebreak
only separately proved regularizable loci may be compared

\medskip
\textbf{Pair of realizations}\par\nopagebreak
local BFJ mixed measure / swept CLD boundary product\par\nopagebreak

\textbf{Status}\par\nopagebreak
equality after Gauss-shell realization and pushforward\par\nopagebreak

\textbf{Exact common locus}\par\nopagebreak
smooth complex germ; fixed normalization ideal and shell; ideal and
algebraizable rational model data, together with the stated joint bounded
closure

\medskip
\textbf{Pair of realizations}\par\nopagebreak
polarized affine-cone model measure / CLD Bedford--Taylor product\par\nopagebreak

\textbf{Status}\par\nopagebreak
literal equality by definition and CLD approximation\par\nopagebreak

\textbf{Exact common locus}\par\nopagebreak
normal affine cone over a trivially valued algebraically closed
characteristic-zero field; good torus action; fixed Reeb field; homogeneous
Fubini--Study data and their stated continuous closure

\medskip
\textbf{Pair of realizations}\par\nopagebreak
polarized affine-cone finite energy / global BJ finite energy\par\nopagebreak

\textbf{Status}\par\nopagebreak
recovery on the quasi-regular branch\par\nopagebreak

\textbf{Exact common locus}\par\nopagebreak
rational Reeb field and the associated polarized orbifold quotient; the
comparison is not asserted for any pair of nonproportional polarizations.

\medskip
\textbf{Pair of realizations}\par\nopagebreak
integrally closed test configurations / rational PL metrics\par\nopagebreak

\textbf{Status}\par\nopagebreak
equivalence modulo domination\par\nopagebreak

\textbf{Exact common locus}\par\nopagebreak
fixed generic fibre, punctured trivialization, polarization, and
linearization

\medskip
\textbf{Pair of realizations}\par\nopagebreak
DXZ test-curve metrics / BJ psh metrics\par\nopagebreak

\textbf{Status}\par\nopagebreak
order isomorphism\par\nopagebreak

\textbf{Exact common locus}\par\nopagebreak
connected projective complex manifold and pseudo-effective
\(\mathbf Q\)-line bundle

\medskip
\textbf{Pair of realizations}\par\nopagebreak
MP tropical spectrum / Berkovich analytification\par\nopagebreak

\textbf{Status}\par\nopagebreak
homeomorphism\par\nopagebreak

\textbf{Exact common locus}\par\nopagebreak
underlying space of a proper complex algebraic variety

\medskip
\textbf{Pair of realizations}\par\nopagebreak
DXZ psh / MP psh\par\nopagebreak

\textbf{Status}\par\nopagebreak
no equivalence asserted\par\nopagebreak

\textbf{Exact common locus}\par\nopagebreak
common underlying space is insufficient

\medskip
\textbf{Pair of realizations}\par\nopagebreak
DXZ finite-energy metrics / MP finite-energy psh functions\par\nopagebreak

\textbf{Status}\par\nopagebreak
bijection with matching radial energy\par\nopagebreak

\textbf{Exact common locus}\par\nopagebreak
K\"ahler class on a compact K\"ahler manifold; finite-energy maximal radial
data; equality is initially characterized on rational divisorial points

\medskip
\textbf{Pair of realizations}\par\nopagebreak
test curves / weak geodesic rays\par\nopagebreak

\textbf{Status}\par\nopagebreak
Legendre equivalence\par\nopagebreak

\textbf{Exact common locus}\par\nopagebreak
maximal/model test curves and the ray class in the RWN theorem

\medskip
\textbf{Pair of realizations}\par\nopagebreak
BFJ / BJ\par\nopagebreak

\textbf{Status}\par\nopagebreak
analogy\par\nopagebreak

\textbf{Exact common locus}\par\nopagebreak
different ground norms and different radial geometry

\endgroup

\section{Two diagnostic examples}

\begin{example}[A two-component model curve]
\label{ex:dict-two-component}
Let \(X\) be a smooth projective curve over a discretely valued field and
let a regular semistable model have special fibre \(E_0+E_1\).
For a semipositive model line bundle \(\mathcal L\), all three model
realizations give
\[
 \deg(\mathcal L|_{E_0})\delta_{x_0}
 +\deg(\mathcal L|_{E_1})\delta_{x_1}.
\tag{63.16}\label{eq:dict-two-component-measure}
\]
The total mass is \(\deg L\).
The equality is literal because every object comes from the same model.
If one now takes an arbitrary singular BFJ limit, the example no longer
proves that a CLD singular wedge exists.
\end{example}

\begin{example}[The product action on \(\mathbf P^1\)]
\label{ex:dict-product-action}
Let \(\mathbf G_m\) act on \((\mathbf P^1,\mathcal O(1))\) with weights
\(0,1\).
The associated filtration has degree-\(m\) weights
\(0,1,\ldots,m\).
Its FS metric, maximal test curve, and product geodesic ray are linked by
\eqref{eq:dict-degeneration-diagram}.
The limiting probability DH measure is Lebesgue measure on \([0,1]\), while
the unnormalized measure has mass
\(V=(\mathcal O(1)^1)=1\), so the two happen to coincide.
For \(\mathcal O(d)\), the coincidence disappears unless the mass factor
\(d\) is displayed.
\end{example}

\section{Forbidden inferences}

We finish with four rules that will be used in the comparison chapters.

\begin{enumerate}
\item Equality of model measures does not imply equality of all singular
products.
The passage from model to singular objects requires a uniqueness theorem for
the chosen limiting procedure.

\item A homeomorphism of analytifications does not identify psh cones.
Positivity depends on a class, approximation family, and envelope operation.

\item Similar intersection formulas over trivial and nontrivial ground
norms do not define a specialization functor between BFJ and BJ theories.
The normalizations of divisorial valuations and the radial action differ.

\item Equality of asymptotic slopes along algebraic rays does not identify
all finite-energy ray spaces.
The ray class, reference normalization, and published extension theorem must
be stated.
\end{enumerate}

\section*{Source note}

The formal intersection formula is
\cite[Theorem~18.5.5 and Proposition~18.5.8]{CLD25}, and the original
model measure is \cite{CL06}.
The BFJ skeletal and Monge--Amp\`ere theories are
\cite{BFJ16,BFJ15}; the common continuous semipositive locus is also treated
in \cite{BE21} under its stated hypotheses.
The normalized local link, nef b-divisors, and local intersection measure are
\cite{BFJ08}; their Gauss-shell comparison is used only in the verified
model and joint bounded domains stated in \Cref{chap:local-valuative-psh}.
The trivially valued FS and psh dictionary is \cite{BJ22}.
The projective complex comparison is \cite{DXZ23}, and the tropical-spectrum
foundations and underlying-space recovery are \cite{MP24}.
The test-curve and ray correspondence is \cite{RWN14}.
Each citation supports only the arrow to which it is attached above.
